\section{Conclusion}

Interstate 2.0 is the evidence-based answer to a question the United States has not formally asked in sixty years: which highways should we invest in next, and why? The ROUTE corpus provides the data-driven prioritization that replaces the political and geographic allocation mechanisms that have governed federal highway spending since the original interstate appropriations.

\subsection{The Portfolio in Brief}

The six-component investment portfolio costs \$253 billion over 30 years --- \$8.4 billion per year. For context, the Infrastructure Investment and Jobs Act of 2021 allocated approximately \$57 billion per year for surface transportation over its five-year authorization period \citep{IIJA2021}. The I2.0 program absorbs less than 15\% of current annual federal highway spending per year, spread over a 30-year horizon rather than concentrated in a five-year appropriation. It is a realignment of existing spending priorities, not a call for new spending authority.

The aggregate NPV at 7\% discount is \$298 billion against \$253 billion in capital: a B/C ratio of 2.2:1. No component in the portfolio falls below 1.2:1. The highest-return components --- diamond interchange upgrades (2.76:1) and compound exposure hardening (3.1:1) --- are also the lowest-capital components, which means their returns can be captured early in the phasing sequence without requiring the full 30-year program commitment.

\subsection{What I2.0 Produces}

A fully implemented I2.0 program produces, by the 2055 completion horizon:
\begin{itemize}
  \item National commercial vehicle throughput increases from 4.8M to 7.2M commercial vehicles per day (50\% increase)
  \item Transcontinental transit time reduces from 4.5 days to 3.5 days on the NY--LA corridor
  \item PTI variance across T1 corridors falls from std dev 0.42 to 0.18 (57\% reduction)
  \item All 15 T1/T1 intersections upgraded from k=1 to k$\geq$2 connectivity
  \item Gulf Coast I-10 and six other compound exposure corridors hardened to 2050 climate standards
  \item Five missing links completed to interstate standard, closing the most consequential structural gaps in the national freight network
  \item 12.4 million transit-dependent Americans brought within 30 miles of an intercity transit connection
\end{itemize}

\subsection{Limitations}

The ROUTE analysis has three principal limitations that should constrain confidence in the specific numbers.

First, the max-flow model uses a single-commodity formulation. Real freight flows are multi-commodity, with different commodities having different value-of-time parameters, different origin-destination structures, and different sensitivity to managed lane tolls. A multi-commodity flow model would produce more precise B/C estimates, particularly for missing link corridors where the commodity mix on the new route is uncertain.

Second, the PTI estimates use the Bureau of Public Roads volume-delay function \citep{BPR1964}, which is a well-validated but simplified model. FHWA's Freight Performance Measures program produces empirical PTI values from probe vehicle data, but probe data coverage is incomplete for rural T1 segments \citep{FHWA_freight2023}. Where FHWA FPM data is available, it was used; where it is not, the BPR function was calibrated to HPMS volume data. The BPR-calibrated estimates may over- or under-estimate PTI on specific segments.

Third, the 2050 climate projections used in the compound exposure analysis carry substantial uncertainty. The NOAA 2022 Sea Level Rise Technical Report provides a range of scenarios rather than a single projection \citep{NOAA_SLR2022}. The D1 scores used in the hardening investment case are based on the median NOAA scenario; under the high-emissions scenario, the Gulf Coast I-10 D1 score reaches 10.0 (maximum) before 2040, which would accelerate the hardening investment timeline.

\subsection{Call to Action}

The United States chose a highway system. The Eisenhower interstate network was a deliberate national choice, executed over thirty years, that remade the economic geography of the country \citep{Michaels2008,TRB_NHS2018}. ROUTE does not argue that the choice was wrong or that a different modal mix would have been better. It argues that the country, having made the choice, should execute it well --- which means investing in the corridors that actually carry freight, closing the gaps the original network left unfilled, and hardening the nodes that the next fifty years of climate exposure will test.

The data supports a specific list of investments, in a specific priority order, with specific benefit-cost estimates. That is what the ROUTE corpus is for. Political and geographic allocation will always compete with evidence-based prioritization for federal highway dollars. The contribution of this program is to make the evidence case rigorous enough that the political case must argue against something specific --- not against an abstraction.
